%==============================================================
\section*{Transparenzerklärung \& methodologische Urheberschaft (KI-Einsatz)}
%==============================================================

Der Ontophänomenalismus (OP) ist das Ergebnis einer jahrzehntelangen philosophischen
und epistemologischen Auseinandersetzung des Autors (Manuel Krüger), die bereits in
der Kindheit mit grundlegenden Fragen nach Existenz, Bewusstsein und der Struktur der
Wirklichkeit begann und sich über viele Jahrzehnte zu einem eigenständigen
Forschungsprogramm entwickelte.

Die grundlegende Konzeption des Ontophänomenalismus, seine drei Grundaxiome, die
Architektur der Dokumentenfamilie sowie die zentralen philosophischen Leitideen wurden
vom Autor entwickelt und fortlaufend ausgearbeitet.

Zur kritischen Prüfung, logischen Konsolidierung, sprachlichen Präzisierung und
technischen Ausarbeitung der LaTeX-Dokumentation werden in diesem Projekt gezielt
große Sprachmodelle (LLMs), darunter Claude, ChatGPT, Gemini, Grok, DeepSeek und
weitere Systeme, eingesetzt.

Die Zusammenarbeit folgt einer klar definierten methodischen Rollenverteilung:

\begin{itemize}
    \item \textbf{Konzeptionelle Leitung (Autor):} Die inhaltliche Verantwortung sowie
    alle grundlegenden konzeptionellen Entscheidungen liegen beim Autor.

    \item \textbf{Adversarial Reviewer (KI):} Die Sprachmodelle werden bewusst nicht als
    Bestätigungsinstanzen eingesetzt, sondern als kritische Gegenposition. Ihre Aufgabe
    besteht darin, logische Schwächen, unbegründete Annahmen, Zirkelschlüsse,
    Inkonsistenzen und alternative Interpretationen aufzudecken. Hierzu werden
    Fragestellungen regelmäßig in neuen, voneinander unabhängigen Dialogen eingebracht,
    um möglichst unvoreingenommene Gegenargumente und Einwände zu erhalten. Vorschläge,
    Kritik und formale Beiträge der Modelle werden ausschließlich nach Prüfung und
    ausdrücklicher Entscheidung des Autors übernommen oder verworfen.

    \item \textbf{Formale Ausarbeitung (KI):} Die sprachliche Präzisierung, logische
    Konsolidierung und LaTeX-Umsetzung erfolgen in enger Zusammenarbeit zwischen Autor
    und KI, um eine möglichst konsistente, präzise und nachvollziehbare Dokumentation
    zu gewährleisten.
\end{itemize}

Dieses Projekt versteht sich als ein Beispiel KI-assistierter menschlicher
Urheberschaft. Die KI-Modelle entwickeln nicht eigenständig die philosophische
Gesamtkonzeption, sondern dienen als methodische Reflexionsinstanzen und
``Adversarial Reviewer`` innerhalb eines bewusst kritisch angelegten
Entwicklungsprozesses. Die philosophische Gesamtverantwortung sowie die endgültige
inhaltliche Gestalt des Ontophänomenalismus verbleiben beim Autor.
